% ================================================================
\section{Background}\label{sec:background}

\subsection{Mathematical Preliminaries}\label{sec:prelim}

This section introduces the three mathematical tools that
the paper builds on: ADMM, the discrete algebraic Riccati
equation, and the log-barrier interior-point method.
Readers familiar with these topics may skip to
Section~\ref{sec:tinympc_background}.

\textbf{ADMM (Alternating Direction Method of Multipliers).}
Many MPC problems take the form
\begin{equation}\label{eq:admm_generic}
  \min_{z} \; f(z) \quad \text{s.t.} \quad z \in \mathcal{C},
\end{equation}
where $f$ is a smooth (often quadratic) cost and $\mathcal{C}$
encodes constraints (e.g., actuator limits, state bounds).
ADMM~\cite{boyd2011} splits this into two alternating steps
by introducing an auxiliary variable $\tilde{z}$ and a dual
variable $\lambda$:
\begin{enumerate}[label=(\alph*)]
  \item \emph{Minimize} $f$ without constraints (smooth, often
    admits a closed-form or Riccati-based solution):
    $z^{k+1} = \arg\min_z \; f(z) + \tfrac{\rho}{2}\|z - \tilde{z}^k + \lambda^k\|^2$.
  \item \emph{Project} onto constraints (often a simple clipping
    operation):
    $\tilde{z}^{k+1} = \mathrm{proj}_{\mathcal{C}}(z^{k+1} + \lambda^k)$.
  \item \emph{Update} the dual variable:
    $\lambda^{k+1} = \lambda^k + z^{k+1} - \tilde{z}^{k+1}$.
\end{enumerate}
The penalty $\rho > 0$ controls the coupling strength between
the smooth and constraint steps.
ADMM converges at a \emph{linear} rate: the primal residual
$\|z - \tilde{z}\|$ decreases geometrically with the number
of iterations~\cite{boyd2011}.
For MPC, the key property is that step~(a) has the structure
of an unconstrained LQR problem, solvable by a Riccati
recursion, while step~(b) is a componentwise clipping
operation.
This separation is what enables the DARE-caching strategy
of TinyMPC.

\textbf{Discrete Algebraic Riccati Equation (DARE).}
Consider the infinite-horizon LQR problem for the linear
system $x_{k+1} = A x_k + B u_k$ with stage cost
$\ell(x,u) = \frac{1}{2}(x^\top Q x + u^\top R u)$,
where $Q \succeq 0$ and $R \succ 0$.
The optimal cost-to-go from state $x$ is
$V^*(x) = \frac{1}{2} x^\top P x$,
where $P$ satisfies the DARE:
\begin{equation}\label{eq:dare_background}
  P = Q + A^\top P A
  - A^\top P B (R + B^\top P B)^{-1} B^\top P A.
\end{equation}
The optimal feedback gain is
$K = (R + B^\top P B)^{-1} B^\top P A$,
so that $u^* = -K x$ minimizes the infinite-horizon cost.
The closed-loop matrix $A_{\mathrm{cl}} = A - BK$ is stable:
all eigenvalues lie strictly inside the unit circle,
guaranteeing exponential convergence
$\|x_k\| \leq c \, \bar\rho(A_{\mathrm{cl}})^k \|x_0\|$
where $\bar\rho(A_{\mathrm{cl}}) < 1$ is the spectral radius.

The DARE is solved offline by iterating the Riccati recursion
$P_{k} = Q + A^\top P_{k+1} A - A^\top P_{k+1} B
(R + B^\top P_{k+1} B)^{-1} B^\top P_{k+1} A$
backward from a terminal condition $P_N$ until convergence.
The converged $P$ and $K$ encode the entire infinite-horizon
optimal controller in two matrices.

\textbf{Log-Barrier Interior-Point Method.}
To handle inequality constraints
$g_j(u) \leq 0$, $j = 1, \ldots, p$,
the log-barrier method~\cite{boyd2004} replaces them with
a smooth penalty:
\begin{equation}\label{eq:barrier_background}
  \min_u \; f(u) - \mu \sum_{j=1}^{p} \log(-g_j(u)),
\end{equation}
where $\mu > 0$ is the barrier parameter.
The logarithmic term is finite only when all constraints
are strictly satisfied ($g_j < 0$), and tends to $+\infty$
as any constraint approaches its boundary ($g_j \to 0^-$).
As $\mu \to 0$, the solution of~\eqref{eq:barrier_background}
approaches the solution of the original constrained problem.

For box constraints $u_{\min} \leq u \leq u_{\max}$,
$g_j(u) = u_j - u_{\max,j}$ and
$g_{j+m}(u) = u_{\min,j} - u_j$,
giving the barrier
$-\mu[\log(u_j - u_{\min,j}) + \log(u_{\max,j} - u_j)]$.
The Hessian (second derivative) of this barrier is a
diagonal matrix whose entries are large near the constraint
boundaries and small far from them.

Each subproblem~\eqref{eq:barrier_background} is smooth
and unconstrained (in the strict interior), so it can be
solved by Newton's method.
Newton's method converges \emph{quadratically}:
the error decreases as $\|e_{k+1}\| \leq C \|e_k\|^2$,
meaning that 2--3 iterations typically reduce a moderate
initial error to machine precision~\cite{boyd2004}.
This is in contrast to ADMM's linear
convergence, where a fixed fraction of the error is removed
per iteration.

\subsection{Related Work}\label{sec:related}

\textbf{Gain scheduling.}
The classical approach to controlling nonlinear systems around
varying operating points is gain scheduling~\cite{rugh2000}:
solve the linear control problem (e.g., LQR or $H_\infty$) at
a grid of operating points offline, then interpolate the resulting
gains online.
Rugh and Shamma~\cite{rugh2000} survey the extensive theory,
including conditions under which interpolated controllers
inherit stability from the frozen-parameter designs.
Our Flat Template MPC is a gain-scheduled ADMM-MPC controller
where the scheduling variable is the flat parameter $\balpha$,
and our Analytic Template is a first-order gain schedule with
an explicit interpolation formula~\eqref{eq:K_interp}.

\textbf{LPV control.}
Linear parameter-varying (LPV) systems generalize gain scheduling
by treating the varying linearization as a parameter-dependent
plant $\dot{x} = A(\theta) x + B(\theta) u$ where $\theta(t)$
is measurable in real time~\cite{shamma1991}.
Our framework fits this structure with $\theta = \balpha$,
but we exploit flatness to guarantee that $\balpha$ is a
\emph{four-dimensional sufficient statistic} for the
equilibrium---rather than requiring the full $n$-dimensional
state as scheduling variable.

\textbf{Riccati sensitivity analysis.}
The dependence of the DARE solution on system parameters has been
studied since Hewer~\cite{hewer1971}, with explicit perturbation
bounds developed by Konstantinov et al.~\cite{konstantinov1993}.
Our gain sensitivities $\partial K / \partial \alpha_i$
in~\eqref{eq:dK} are a numerical instance of this theory,
approximated by central differences.
The key practical insight is that for quadrotor dynamics,
$K(\balpha)$ varies smoothly and slowly with $\balpha$,
so a first-order expansion suffices (4\% mean relative error
over the tested range).

\textbf{Flatness-based tracking control.}
Differential flatness has been widely used for quadrotor
trajectory generation~\cite{mellinger2011,mueller2015} and
feedforward computation~\cite{faessler2018}.
Mellinger and Kumar~\cite{mellinger2011} use the flat map to
compute reference states and inputs, then track with a fixed-gain
controller.
Our approach extends this by additionally using the flat map
to parameterize the \emph{feedback gain itself},
adapting $K$ to the current operating condition through
the DARE sensitivity.

\textbf{Explicit MPC.}
An alternative to online optimization is explicit MPC~\cite{bemporad2002},
which precomputes the optimal control law as a piecewise-affine
function of the state via multi-parametric programming.
The Flat Template MPC shares the ``precompute offline, look up online''
philosophy, but avoids the exponential complexity of explicit MPC
by storing only DARE caches (not full polyhedral partitions)
and running a small number of ADMM iterations online.

\subsection{TinyMPC}\label{sec:tinympc_background}

TinyMPC~\cite{tinympc} solves the constrained linear MPC problem,
formulated as a quadratic program (QP)~\cite{osqp,boyd2011},
\begin{align}
  \min_{x_{1:N}, u_{1:N-1}} \;\;
  & \tfrac{1}{2} x_N^\top Q_f x_N
    + \sum_{k=1}^{N-1} \bigl(
      \tfrac{1}{2} x_k^\top Q x_k
      + \tfrac{1}{2} u_k^\top R u_k \bigr) \notag \\
  \text{s.t.} \;\;
  & x_{k+1} = A x_k + B u_k, \notag \\
  & x_k \in \mathcal{X},\;\; u_k \in \mathcal{U},
  \label{eq:mpc}
\end{align}
via ADMM, where $\rho > 0$ is the ADMM penalty parameter.
The key insight is that, for a sufficiently long horizon $N$,
the backward Riccati recursion converges to a steady-state
solution $(K_\infty, P_\infty)$.
The matrix $P_\infty$ is the infinite-horizon optimal
cost-to-go: it compresses the effect of all future stages
into a single $n \times n$ matrix, so that
$\frac{1}{2} \delta x^\top P_\infty \, \delta x$ approximates
the minimum cost from state $\delta x$ to the origin
over an infinite horizon.
TinyMPC solves the discrete algebraic Riccati equation (DARE)
for $P_\infty$ \emph{offline} and caches the derived quantities:
\begin{align}
  K_\infty &= (R_\rho + B^\top P_\infty B)^{-1} B^\top P_\infty A, \label{eq:kinf} \\
  C_1      &= (R_\rho + B^\top P_\infty B)^{-1}, \notag \\
  C_2      &= (A - B K_\infty)^\top, \notag
\end{align}
where $R_\rho = R + \rho I$ is the ADMM-augmented input cost
and $Q_\rho = Q + \rho I$ is the ADMM-augmented state cost.
The gain $K_\infty$ is the optimal LQR gain for the augmented
problem; $C_1$ and $C_2$ are auxiliary matrices that eliminate
all matrix inversions from the online computation.
Each ADMM iteration then requires only matrix--vector products:
\begin{enumerate}[label=(\roman*)]
  \item \emph{Backward pass}: compute feedforward corrections
    $d_k = C_1(B^\top p_{k+1} + r_k)$,
    $p_k = q_k + C_2 \, p_{k+1} - K_\infty^\top r_k$.
  \item \emph{Forward pass}: $u_k = -K_\infty x_k - d_k$,
    $x_{k+1} = A x_k + B u_k$.
  \item \emph{Slack update}: project onto constraints.
  \item \emph{Dual update}: gradient ascent on multipliers.
\end{enumerate}

The entire online computation is \emph{matrix-inversion free}.
However, the matrices $A, B$ (and hence all cached quantities)
are fixed at a single linearization point.

\subsection{Differential Flatness of Quadrotor Dynamics}\label{sec:flatness}

A quadrotor with state
$(p, \dot{p}, R, \omega) \in \R^3 \times \R^3 \times \mathrm{SO}(3) \times \R^3$
and collective thrust $T$ along the body $z$-axis satisfies
\begin{equation}\label{eq:trans}
  m \ddot{p} = T R e_3 - m g e_3.
\end{equation}
The system is differentially flat with flat output
$\by = (p_x, p_y, p_z, \psi)$~\cite{mellinger2011,murray1995}.
This means that every state and input can be expressed as an
algebraic function of $\by$ and a finite number of its time derivatives.

The flat map proceeds as follows.
Define the \emph{total acceleration vector}
\begin{equation}\label{eq:tvec}
  \bt = \ddot{p} + g e_3 = \begin{pmatrix} a_x \\ a_y \\ a_z + g \end{pmatrix},
\end{equation}
where $(a_x, a_y, a_z)$ is the commanded translational acceleration
in the world frame.
Then:
\begin{enumerate}
  \item \textbf{Thrust:} $T = m \| \bt \|$.
  \item \textbf{Attitude:} The body $z$-axis is
    $z_B = \bt / \|\bt\|$.
    Given yaw $\psi$, the rotation $R$ is fully determined.
    In Euler angles (ZYX convention):
    \begin{align}
      \phi^* &= \arcsin\!\bigl(
        (a_x \sin\psi - a_y \cos\psi) / \|\bt\| \bigr), \label{eq:phi_eq} \\
      \theta^* &= \arctan2\!\bigl(
        a_x \cos\psi + a_y \sin\psi, \; a_z + g \bigr). \label{eq:theta_eq}
    \end{align}
  \item \textbf{Angular velocity and torques:}
    determined by $p^{(3)}$ and $p^{(4)}$ (jerk and snap) through the
    kinematic and Euler equations; at equilibrium, $\omega^* = 0$.
\end{enumerate}

We define the \emph{flat parameter}
$\balpha = (a_x, a_y, a_z, \psi) \in \R^4$,
which fully specifies the equilibrium state and input:
\begin{equation}\label{eq:eq_map}
  \Phi: \balpha \;\mapsto\; (x^*(\balpha),\; u^*(\balpha)).
\end{equation}
